To determine how the spread of selling depended on the number of
interacting traders, we compared grid sizes with technical traders held
and found a sharper majority response on larger grids
(Fig.~\ref{fig:finite-size}). We refined the fractions near each
response while retaining the same starting markets. At the original
3,000-step horizon, the interpolated fraction at which half the markets
reached a majority, $f_{50}$, changed from 0.216 at $L=32$ to
0.166 at $L=128$. The width separating majority probabilities
of 0.25 and 0.75 changed from 0.023 [0.018, 0.027] to
0.008 [0.004, 0.010], with brackets giving 95\% intervals from
resampling whole starting markets (Fig.~\ref{fig:finite-size}A,C;
Supplementary Table~\ref{tab:finite-size-scaling}). Variation in the
final mean Sell fraction also changed: its variance multiplied by the
number of interior traders, the susceptibility, had sampled maxima of
21.2 and 241.2 at the
endpoint sizes (Fig.~\ref{fig:finite-size}B,D). Across all size and
feedback conditions, the largest absolute difference between logistic
and interpolated midpoints was 0.0035, indicating the sensitivity
of the estimate to the imposed curve shape
(Supplementary Tables~\ref{tab:finite-size-scaling} and~\ref{tab:revision-fits}). Observation time
remained consequential. At $L=64$, extending the held markets from
3,000 to 9,000 steps changed the majority midpoint from
0.178 to 0.167; alternative
Sell thresholds and size-dependent horizons were measured separately
(Supplementary Table~\ref{tab:revision-horizons}). The size response
therefore described a crossover within the observation window, while
the continued movement with time prevented its interpretation as a
fixed threshold for indefinitely long runs.
